\honest{Zut Allais!}{Eliezer Yudkowsky}{2008}

I did not expect yesterday's commenters to defend the Allais choices. It helps to know two things from the study of heuristics and biases: subjects defend incoherent preferences ``even when they're really silly,'' and people put very high values on small shifts in probability away from 0 or 1. \nb{For the Allais pattern the evidence on persistence is mixed: some studies found that violations survive arguments for the axiom, one found a tendency to correct them. Some researchers read persistence as support for Allais.}

First, incoherence. People choose one bet but put a higher price on another, because different features stand out when you ask ``Which do you prefer?'' and ``How much would you pay?'' My books are packed, so I give an example from memory, twice marked ``IIRC'': a 1/3 chance to win \$18 and 2/3 to lose \$1.50, against 19/20 to win \$4 and 1/20 to lose \$0.25. People would rather play the second but price the first higher. \nb{Kahneman and Tversky credit this pattern to Lichtenstein and Slovic, and say their own theory does not treat bidding. The post's two bets could not be traced.}

Sell such a person the first bet at their price, swap it for the second, buy that back at their lower price, and repeat: a money pump. As Steve Omohundro put it, if you prefer Oakland to San Francisco, San Jose to Oakland and San Francisco to San Jose, you will spend a lot on taxis.

Some subjects give up the pattern when this is explained; some defend it. In a transcript from Las Vegas, where gamblers played such bets for real money on a roulette wheel, a subject admits ``It shows my reasoning process isn't so good'' but has ``no qualms,'' doubts that the experimenter could persuade him the pattern is irrational, and, walked through the pump, pays 550 points for one bet and sells the other back for 401, leaving the experimenter ``ahead 149 points.'' ``That's good reasoning on my part,'' he laughs. You want to scream, ``Just give up already! Intuition isn't always right!''

Second, certainty. ``I believe'' one experiment found that a shift from 100\% to 99\% weighed more than a shift from 80\% to 20\%. \nb{The nearest published finding compares a drop from 1.0 to .25 with one from .8 to .2 (Kahneman and Tversky, 1979). Nothing close to the post's version could be found.} The trouble is that ``one time's certainty is another time's probability.'' In yesterday's pump you pay to switch to B, preferring 33 per cent of \$27,000 to 34 per cent of \$24,000; then the die removes the 66 per cent chance of nothing that both options shared, A becomes glorious certainty, and you pay to switch back. Told the plan in advance, would you ``prefer to reconsider''? Valuing a shift from 24 to 23 per cent less than one from 100 to 99 per cent, giving extra weight near the ends of the scale, lets me remove probability a bit at a time until your preference flips.

You are a flawed piece of machinery, and your intuitions are not direct information about good choices. If you believe otherwise, I have some gambling games to play with you. A sure \$400, or 80\% of \$500 and 20\% of \$300? Most take the \$400. Start \$500 richer and choose between losing \$100 for sure or a 20\% chance of losing \$200? Most gamble. The outcomes are the same. \nb{Kahneman and Tversky tested this design in 1979 with the result the post describes, using other numbers.}

So multiply utilities by probabilities; do not be embarrassed to use clean math. For the Allais problem, ask whether one unit of the difference between \$24,000 and nothing outweighs 33 units of the difference between \$24,000 and \$27,000, and choose the same way in both cases. Use logarithmic utility of total assets, not of changes in assets, or you will be inconsistent again: with plenty of money, pick B. ``If \$24,000 would double your existing assets, pick A.'' \nb{Log utility gives B at those assets too. It favors A only for assets under about \$550.}

Commenters wrote ``U(Certainty)'' into an equation. But utility and expected utility are different types: utilities attach to particular states of the world. ``You cannot feed a probability of 1 into a utility function.'' And the price of leaving the Bayesian Way here was paying to throw a switch and throw it back. \nb{Published defenders of the Allais choices put the value of certainty into the outcomes instead, for example the disappointment of losing a sure thing. The post answers the version about feelings, next.}

Paying for a warm feeling of certainty is human; it depends whether you care more about satisfying your intuitions or achieving the goal. Gambling in Las Vegas for fun, do not bother with expected utility. You will lose anyway. But if 24,000 lives were at stake, the certainty effect is ``even stronger.'' Will you pay one life to throw the switch and another to switch it back? \nb{No source is given. In Tversky and Kahneman's 1981 study, people preferred the sure option when lives were described as saved and the gamble when they were described as deaths.}

Preferences that reverse are like driving from San Jose to San Francisco to Oakland and back, forever: warm feelings, but no destination. You are not steering the future, just running in circles. When a life is at stake, your intuitions are a greasy lens. There are mathematical laws for steering the future efficiently, and when something more important than your feelings is at stake, you should care about them.
